\documentclass[12pt,a4paper]{article}
\usepackage[UTF8]{ctex}
\usepackage{amsmath,amssymb,amsthm}
\usepackage{geometry}
\usepackage{graphicx}

\geometry{margin=2.5cm}

\title{非周期信号中相位差导致的滞后}
\author{第11章补充说明}
\date{}

\begin{document}

\maketitle

\section{问题提出}

对于周期信号，我们知道了相位差导致的时域滞后：
\begin{equation}
\Delta t = \frac{|\Delta \phi|}{\omega} = \frac{|\Delta \phi| \cdot T}{2\pi}
\end{equation}

但对于\textbf{非周期信号}，情况更复杂。本文将详细分析非周期信号中相位差的影响。

\section{非周期信号的频率分析}

\subsection{傅里叶变换}

非周期信号可以通过傅里叶变换分解为连续频率分量：
\begin{equation}
x(t) = \frac{1}{2\pi} \int_{-\infty}^{\infty} X(j\omega) e^{j\omega t} d\omega
\end{equation}

其中$X(j\omega)$是信号的频谱。

\subsection{系统响应}

对于线性时不变系统$G(j\omega)$，输出为：
\begin{equation}
y(t) = \frac{1}{2\pi} \int_{-\infty}^{\infty} G(j\omega) X(j\omega) e^{j\omega t} d\omega
\end{equation}

\textbf{关键观察}：
\begin{itemize}
\item 每个频率分量$\omega$有不同的相位滞后$\angle G(j\omega)$
\item 不同频率分量的滞后不同
\item 输出信号是所有这些分量的叠加
\end{itemize}

\section{积分器对非周期信号的影响}

\subsection{积分器的频率响应}

对于积分器$G(s) = 1/s$：
\begin{equation}
G(j\omega) = \frac{1}{j\omega} = \frac{1}{\omega} e^{-j\pi/2}
\end{equation}

\textbf{特性}：
\begin{itemize}
\item 相位滞后：$-90^\circ$（对所有频率）
\item 幅值：$1/\omega$（随频率变化）
\end{itemize}

\subsection{不同频率分量的滞后}

对于频率为$\omega$的分量：
\begin{equation}
\Delta t(\omega) = \frac{\pi/2}{\omega} = \frac{T(\omega)}{4}
\end{equation}

其中$T(\omega) = 2\pi/\omega$是该频率分量的周期。

\textbf{关键问题}：
\begin{itemize}
\item 不同频率分量有不同的滞后
\item 低频分量：滞后大（因为周期大）
\item 高频分量：滞后小（因为周期小）
\end{itemize}

\section{典型非周期信号的分析}

\subsection{阶跃信号}

\textbf{输入}：单位阶跃信号
\begin{equation}
x(t) = \begin{cases}
0 & t < 0 \\
1 & t \geq 0
\end{cases}
\end{equation}

\textbf{频谱}：
\begin{equation}
X(j\omega) = \pi \delta(\omega) + \frac{1}{j\omega}
\end{equation}

\textbf{积分器输出}：
\begin{equation}
y(t) = \int_0^t x(\tau) d\tau = \begin{cases}
0 & t < 0 \\
t & t \geq 0
\end{cases}
\end{equation}

\textbf{分析}：
\begin{itemize}
\item 阶跃信号包含所有频率分量（从0到$\infty$）
\item 每个频率分量有不同的滞后
\item 但输出是简单的斜坡函数，没有明显的"滞后"
\item 这是因为阶跃信号的特性：所有频率分量同时开始
\end{itemize}

\subsection{脉冲信号}

\textbf{输入}：单位脉冲（Dirac delta）
\begin{equation}
x(t) = \delta(t)
\end{equation}

\textbf{频谱}：
\begin{equation}
X(j\omega) = 1 \quad \text{（所有频率分量幅值相同）}
\end{equation}

\textbf{积分器输出}：
\begin{equation}
y(t) = \int_{-\infty}^t \delta(\tau) d\tau = u(t) = \begin{cases}
0 & t < 0 \\
1 & t \geq 0
\end{cases}
\end{equation}

\textbf{分析}：
\begin{itemize}
\item 脉冲包含所有频率分量
\item 积分器将脉冲转换为阶跃
\item 没有明显的"滞后"，因为脉冲是瞬时的
\end{itemize}

\subsection{斜坡信号}

\textbf{输入}：斜坡信号
\begin{equation}
x(t) = \begin{cases}
0 & t < 0 \\
t & t \geq 0
\end{cases}
\end{equation}

\textbf{积分器输出}：
\begin{equation}
y(t) = \int_0^t \tau d\tau = \frac{t^2}{2} \quad t \geq 0
\end{equation}

\textbf{分析}：
\begin{itemize}
\item 输出是二次函数
\item 没有明显的"滞后"概念
\end{itemize}

\section{实际应用：导纳控制中的非周期信号}

\subsection{用户施加的力}

在实际应用中，用户施加的力$f_{\text{ext}}(t)$通常是非周期的：
\begin{itemize}
\item 可能是突然的推拉（类似阶跃）
\item 可能是缓慢变化（类似斜坡）
\item 可能是复杂的时变信号
\end{itemize}

\subsection{导纳控制响应}

导纳控制计算：
\begin{equation}
\ddot{x}_d = M^{-1}(f_{\text{ext}} - B\dot{x} - Kx)
\end{equation}

然后积分：
\begin{equation}
\dot{x}_d(t) = \dot{x}_d(0) + \int_0^t \ddot{x}_d(\tau) d\tau
\end{equation}

\subsection{滞后的表现}

对于非周期信号，相位滞后的表现：

\begin{enumerate}
\item \textbf{阶跃力输入}：
\begin{itemize}
\item 用户突然施加力（阶跃）
\item 系统响应有延迟
\item 延迟主要来自系统的动态响应，而不是相位滞后
\end{itemize}

\item \textbf{缓慢变化的力}：
\begin{itemize}
\item 如果力变化很慢（低频），可以近似为周期信号
\item 相位滞后$\Delta t \approx T/4$，其中$T$是"有效周期"
\item 对于非常慢的变化，滞后可能很大
\end{itemize}

\item \textbf{快速变化的力}：
\begin{itemize}
\item 如果力变化很快（高频），滞后较小
\item 但系统可能跟不上快速变化
\end{itemize}
\end{enumerate}

\section{线性相位系统}

\subsection{定义}

如果系统的相位响应是频率的线性函数：
\begin{equation}
\angle G(j\omega) = -\tau \omega
\end{equation}

其中$\tau$是常数，称为\textbf{群延迟}（group delay）。

\subsection{时域滞后}

对于线性相位系统：
\begin{equation}
\Delta t(\omega) = \frac{|\angle G(j\omega)|}{\omega} = \frac{\tau \omega}{\omega} = \tau
\end{equation}

\textbf{关键特性}：
\begin{itemize}
\item 所有频率分量的时域滞后相同：$\Delta t = \tau$
\item 信号只是平移，形状不变
\item 这是理想情况
\end{itemize}

\subsection{积分器不是线性相位系统}

对于积分器$G(s) = 1/s$：
\begin{equation}
\angle G(j\omega) = -\frac{\pi}{2} \quad \text{（恒定，不是$\omega$的线性函数）}
\end{equation}

因此：
\begin{equation}
\Delta t(\omega) = \frac{\pi/2}{\omega} = \frac{T(\omega)}{4}
\end{equation}

\textbf{结论}：
\begin{itemize}
\item 不同频率分量有不同的滞后
\item 低频分量滞后大
\item 高频分量滞后小
\item 信号形状会改变（不仅仅是平移）
\end{itemize}

\section{信号失真的原因}

\subsection{为什么信号形状会改变？}

对于非周期信号，如果不同频率分量有不同的滞后：
\begin{itemize}
\item 低频分量滞后大
\item 高频分量滞后小
\item 输出信号是这些不同滞后分量的叠加
\item 因此信号形状会改变
\end{itemize}

\textbf{例子}：
\begin{itemize}
\item 输入：方波（包含基频和谐波）
\item 积分后：三角波（形状改变）
\item 不同频率分量的滞后不同，导致形状改变
\end{itemize}

\subsection{群延迟 vs. 相位延迟}

\textbf{相位延迟}（Phase Delay）：
\begin{equation}
\tau_p(\omega) = -\frac{\angle G(j\omega)}{\omega}
\end{equation}
每个频率分量的延迟。

\textbf{群延迟}（Group Delay）：
\begin{equation}
\tau_g(\omega) = -\frac{d\angle G(j\omega)}{d\omega}
\end{equation}
信号包络的延迟。

\textbf{对于积分器}：
\begin{itemize}
\item 相位延迟：$\tau_p(\omega) = \frac{\pi/2}{\omega}$（随频率变化）
\item 群延迟：$\tau_g(\omega) = \frac{\pi/2}{\omega}$（也随频率变化）
\item 两者都随频率变化，因此信号会失真
\end{itemize}

\section{实际测量和估计}

\subsection{如何测量滞后？}

对于非周期信号，滞后的定义更复杂：

\begin{enumerate}
\item \textbf{峰值滞后}：
\begin{itemize}
\item 测量输入和输出峰值之间的时间差
\item 适用于有明显峰值的信号
\end{itemize}

\item \textbf{上升时间滞后}：
\begin{itemize}
\item 测量从10\%到90\%上升时间的中点
\item 适用于阶跃响应
\end{itemize}

\item \textbf{相关性分析}：
\begin{itemize}
\item 计算输入和输出的互相关函数
\item 峰值位置对应滞后
\end{itemize}
\end{enumerate}

\subsection{有效滞后}

对于非周期信号，可以定义\textbf{有效滞后}：

\begin{equation}
\Delta t_{\text{eff}} = \frac{\int_0^{\infty} \omega \cdot \Delta t(\omega) \cdot |X(j\omega)|^2 d\omega}{\int_0^{\infty} \omega \cdot |X(j\omega)|^2 d\omega}
\end{equation}

这是加权平均滞后，权重是信号的能量分布。

\section{数值例子}

\subsection{例子1：缓慢变化的力}

假设用户施加的力缓慢变化，可以近似为正弦波：
\begin{equation}
f_{\text{ext}}(t) = F_0 \sin(0.1t) \quad \text{（频率$f = 0.1/(2\pi) \approx 0.016$ Hz）}
\end{equation}

\textbf{周期}：$T = 2\pi/0.1 \approx 62.8$ s

\textbf{相位滞后}：$-90^\circ$

\textbf{时域滞后}：
\begin{equation}
\Delta t = \frac{T}{4} = \frac{62.8}{4} \approx 15.7 \text{ s}
\end{equation}

\textbf{结论}：对于缓慢变化的信号，滞后很大。

\subsection{例子2：快速变化的力}

假设力快速变化：
\begin{equation}
f_{\text{ext}}(t) = F_0 \sin(10t) \quad \text{（频率$f = 10/(2\pi) \approx 1.59$ Hz）}
\end{equation}

\textbf{周期}：$T = 2\pi/10 \approx 0.628$ s

\textbf{时域滞后}：
\begin{equation}
\Delta t = \frac{T}{4} = \frac{0.628}{4} \approx 0.157 \text{ s}
\end{equation}

\textbf{结论}：对于快速变化的信号，滞后较小。

\subsection{例子3：阶跃力}

假设用户突然施加恒力：
\begin{equation}
f_{\text{ext}}(t) = \begin{cases}
0 & t < 0 \\
F_0 & t \geq 0
\end{cases}
\end{equation}

\textbf{分析}：
\begin{itemize}
\item 阶跃包含所有频率分量
\item 低频分量滞后大（如$f = 0.1$ Hz，滞后$\approx 2.5$ s）
\item 高频分量滞后小（如$f = 10$ Hz，滞后$\approx 0.025$ s）
\item 系统响应是这些分量的叠加
\item 实际滞后取决于系统的主要频率成分
\end{itemize}

\section{减少滞后的方法}

\subsection{方法1：预测补偿}

使用预测来补偿滞后：
\begin{equation}
\dot{x}_d(t) = \int_0^t \ddot{x}_d(\tau) d\tau + \alpha \ddot{x}_d(t)
\end{equation}

其中$\alpha$是预测系数。

\subsection{方法2：高通滤波}

对输入信号进行高通滤波，减少低频分量：
\begin{itemize}
\item 低频分量滞后大
\item 减少低频分量可以减少平均滞后
\item 但可能丢失重要信息
\end{itemize}

\subsection{方法3：自适应控制}

根据信号特性调整控制参数：
\begin{itemize}
\item 对于快速变化：使用较小的$M$（减少惯性）
\item 对于缓慢变化：使用较大的$B$（增加阻尼）
\end{itemize}

\section{总结}

\subsection{核心结论}

\begin{enumerate}
\item \textbf{非周期信号的复杂性}：
\begin{itemize}
\item 包含多个频率分量
\item 每个分量有不同的滞后
\item 信号形状可能改变
\end{itemize}

\item \textbf{滞后的估计}：
\begin{itemize}
\item 对于缓慢变化：可以近似为周期信号，滞后$\approx T/4$
\item 对于快速变化：滞后较小
\item 对于阶跃等：滞后取决于主要频率成分
\end{itemize}

\item \textbf{实际影响}：
\begin{itemize}
\item 低频分量滞后大（可能几秒到几十秒）
\item 高频分量滞后小（可能几毫秒）
\item 平均滞后取决于信号的能量分布
\end{itemize}
\end{enumerate}

\subsection{关键公式}

对于频率为$\omega$的分量：
\begin{equation}
\Delta t(\omega) = \frac{\pi/2}{\omega} = \frac{T(\omega)}{4}
\end{equation}

对于非周期信号，有效滞后：
\begin{equation}
\Delta t_{\text{eff}} \approx \frac{T_{\text{dominant}}}{4}
\end{equation}

其中$T_{\text{dominant}}$是信号主要频率分量的周期。

\subsection{实际建议}

\begin{enumerate}
\item \textbf{对于缓慢变化的信号}：
\begin{itemize}
\item 滞后可能很大（秒级）
\item 需要考虑这个影响
\item 可以使用预测补偿
\end{itemize}

\item \textbf{对于快速变化的信号}：
\begin{itemize}
\item 滞后较小（毫秒级）
\item 但系统可能跟不上
\item 需要提高系统带宽
\end{itemize}

\item \textbf{对于阶跃等瞬态信号}：
\begin{itemize}
\item 滞后取决于系统响应
\item 主要受系统动态特性影响
\item 相位滞后是次要因素
\end{itemize}
\end{enumerate}

\end{document}

